% !PRINTED 406
\weberpage{406}
%<<CONT

mehrfach enth\"alt, und dass in Folge dessen auch $R\,(t, \tau)$ durch
kein Quadrat theilbar ist. Angenommen n\"amlich, es h\"atten
\[
\begin{aligned}
\varPhi\,(t,\tau) &= \varphi\,(t)\,\psi\,(\tau) - \varphi\,(\tau)\,\psi\,(t) \\
\varPhi'(t,\tau) &= \varphi'(t)\,\psi\,(\tau) - \varphi\,(\tau)\,\psi'(t)
\end{aligned}
\]
als Functionen von $t$ einen gemeinsamen Theiler $P$, so m\"usste $P$
auch Theiler von
\[
\psi'(t)\,\varPhi - \psi(t)\,\varPhi' = [\varphi\,(t)\,\psi'(t) - \psi\,(t)\,\varphi'(t)]\,\psi\,(\tau)
\]
sein. Es w\"are daher $P$ Theiler von $\varphi\,(t)\,\psi'(t) - \psi\,(t)\,\varphi'(t)$, und
k\"onnte also von $\tau$ unabh\"angig angenommen werden.

Nehmen wir nun zwei Werthe $\tau_1, \tau_2$ von $\tau$ so an, dass
$\varPhi\,(\tau_2, \tau_1)$ von Null verschieden wird, so ist nach (3)
\[
\varPhi\,(t,\tau_1)\,\varphi\,(\tau_2) - \varPhi\,(t,\tau_2)\,\varphi\,(\tau_1) = \varphi\,(t)\,\varPhi\,(\tau_2,\tau_1).
\]
Darin ist die linke Seite durch $P$ theilbar, und also ist
auch $\varphi\,(t)$ und folglich $\psi\,(t)$ durch $P$ theilbar, was der Voraussetzung
 widerspricht, dass diese beiden Functionen ohne gemeinschaftlichen
 Theiler sein sollen.

Daraus folgt beil\"aufig, dass $R\,(t,\tau) = -R\,(\tau,t)$ sein muss,
da $R$ durch $t - \tau$, aber nicht durch $(t-\tau)^2$ theilbar ist.

Um die Function $R\,(t,\tau)$ zu bilden, kann man den gr\"ossten
gemeinschaftlichen Theiler der Functionen $\varphi_h(t) - v_h\,\psi_h(t)$ aufsuchen
, woraus folgt, dass $R\,(t,\tau)$, abgesehen von einem von $t$
unabh\"angigen Factor, rational durch $v, v_1, \ldots, v_{\nu-1}$ dargestellt
werden kann. Sind also $a$ und $b$ irgend zwei feste numerische
Werthe, so ist
\weq{5}{\frac{R\,(a,\tau)}{R\,(b,\tau)} = \vartheta = \chi\,(v, v_1, \ldots, v_{\nu-1})}
eine rationale Function von $v, v_1, \ldots, v_{\nu-1}$. Dabei ist, wenn $\xi$
eine neue Variable bedeutet:
\weq{6}{R\,(a,\xi) - \vartheta\,R\,(b,\xi) = X}
in Bezug auf $\xi$ vom $\mu^{\mathrm{ten}}$ Grade.

Ist $\tau$ ein willk\"urlicher Werth, so m\"ogen die Wurzeln der
Gleichung $R\,(t,\tau) = 0$
\weq{7}{t = \tau, \tau', \tau'', \ldots}
sein, so dass f\"ur jeden Index $h$
\[
\varPhi_h\,(\tau, \tau), \quad \varPhi_h\,(\tau', \tau), \quad \varPhi_h\,(\tau'', \tau), \ldots
\]
verschwinden.
